% !TEX root = ../main.tex
%%%%%%%%%%%%%%%%%%%%%%%%%%%%%%%%%%%%%%%%%%%%%%%%%%%%%%%%%%%%%%%%%%%%%%%
% Relation to the literature
%%%%%%%%%%%%%%%%%%%%%%%%%%%%%%%%%%%%%%%%%%%%%%%%%%%%%%%%%%%%%%%%%%%%%%%

\section{Relation to the literature}
\label{sec:literature}

\paragraph{Agent-based and stock-flow-consistent macroeconomics.}
The model belongs to the tradition of macroeconomic ABMs that impose
stock-flow consistency on a disaggregated economy
\citep{GodleyLavoie2007,Caiani2016}, and shares with the Schumpeter-meeting-%
Keynes family \citep{Dosi2010} the ambition of letting aggregate regularities
emerge from interacting heterogeneous agents rather than from a
representative optimiser. It is deliberately smaller than those models: there
is one good, one price (the numéraire), no credit and no entry or exit. The
purpose is not to reproduce a full macroeconomy but to isolate a single
mechanism --- endogenous capital accumulation inside a demand-driven
economy --- with enough structure to make its sign an empirical question.

Two of those simplifications are omissions relative to the base model in
particular. In \citet{Teglio2025} household rationality --- four graded levels,
from zero-intelligence to semi-rational --- and the symmetry of the interaction
network are the \emph{independent variables} of the study, not implementation
details; here both are held fixed. The choice is deliberate: to attribute a
measured effect to the accumulation channel the interaction topology has to be
held still, or the effect is not attributable to anything. It is not costless,
and the cost is measured internally rather than asserted ---
\cref{sec:heterogeneity} shows the heterogeneity cascade running precisely
through the fixed spending shares, the very object frozen here, so that with the
network held symmetric a failing firm's demand is destroyed rather than
rerouted. The fixed price inherited from the base model calls for the same
genealogical note: there the numéraire is innocuous, because the wage is not a
price but a share of revenues distributed in equal parts to a firm's employees
\citep{Teglio2025}, so ``real equals nominal'' costs nothing when no relative
price moves. A wage curve changes that --- once the wage responds to
unemployment it is itself a relative price, and holding the numéraire fixed is
no longer free --- which is why the wage-flexibility result is reported as
bracketed by the pass-through convention (\cref{sec:frontierlocal}) rather than
as a property of the labour market alone.

\paragraph{The elasticity of substitution.}
Comparing economies that differ in $\sigma$ is not straightforward: in the
textbook CES, varying $\sigma$ at fixed distribution and efficiency parameters
also moves the implied factor shares, so any behavioural difference is
unattributable. We therefore use the \emph{normalised} CES
\citep{KlumpDeLaGrandville2000,KlumpSaam2008,Klump2012}, which fixes a common
base point through which every $\sigma$-variant passes with the same factor
shares. We follow \citet{Temple2012} in reporting, rather than suppressing,
the caveat that normalisation is a comparison device and not an
identification strategy: it makes $\sigma$-variants commensurable, it does not
make $\sigma$ a deep parameter. \Cref{sec:frontier} accordingly reports the
sensitivity of our central estimate to the choice of anchor.

Empirically, the literature places $\sigma$ below unity: \citet{Chirinko2008}
weighs the evidence at $0.40$--$0.60$; \citet{ChirinkoMallick2017} estimate
$\approx 0.40$; the meta-regression of \citet{Knoblach2020}, covering 2{,}419
estimates from 77 studies, gives $0.45$--$0.87$ and rejects the
Cobb-Douglas hypothesis --- a rejection consistent with earlier direct tests
of the aggregate specification \citep{DuffyPapageorgiou2000,Antras2004}. A
dissenting macro literature requires $\sigma > 1$
to explain the global decline of the labour share
\citep{KarabarbounisNeiman2014}. We treat this disagreement as a reason to
\emph{sweep} $\sigma$ rather than choose it.

\paragraph{The wage curve.}
Wages are endogenised through the wage curve of
\citet{BlanchflowerOswald1994}, a \emph{level} relation between the local wage
and local unemployment with an elasticity near $-0.10$ ($\approx -0.07$ in the
meta-analysis of \citealp{NijkampPoot2005}). We use a level relation rather
than a Phillips-type change relation deliberately: without trend inflation, a
change relation produces perpetual wage drift at any $U \neq U^{\ast}$ and has
no steady state, which would make the comparative-statics design of this
paper impossible.

\paragraph{Distribution and effective demand.}
The question the model is built to answer is Kaleckian
\citep{Kalecki1971,BhaduriMarglin1990}: whether a redistribution toward
profits is expansionary or contractionary depends on whether the induced
investment exceeds the lost consumption. Our contribution is not a new theory
of that trade-off but a measurement of where its sign turns inside an
explicit, stock-flow-consistent, disaggregated economy --- and, more
importantly, a measurement of how little of that sign survives when the rest
of the parameter space is allowed to move.
